
\documentclass{article}
\usepackage{colortbl}
\usepackage{makecell}
\usepackage{multirow}
\usepackage{supertabular}

\begin{document}

\newcounter{utterance}

\centering \large Interaction Transcript for game `function\_detective', experiment `string\_passive\_labeled\_pairs\_oneshot', episode 77 with gpt{-}4.1{-}nano.
\vspace{24pt}

{ \footnotesize  \setcounter{utterance}{1}
\setlength{\tabcolsep}{0pt}
\begin{supertabular}{c@{$\;$}|p{.15\linewidth}@{}p{.15\linewidth}p{.15\linewidth}p{.15\linewidth}p{.15\linewidth}p{.15\linewidth}}
    \# & \multicolumn{2}{c}{Player} && \multicolumn{2}{c}{Game Master} \\
    \hline

    \theutterance \stepcounter{utterance}  
    & & & \multicolumn{4}{p{0.6\linewidth}}{
        \cellcolor[rgb]{0.9,0.9,0.9}{
            \makecell[{{p{\linewidth}}}]{
                \texttt{\tiny{[P1$\langle$GM]}}
                \texttt{I have a hidden Python function for you to figure out.} \\
\\ 
\texttt{Your job is to infer what it does and then submit a correct Python implementation.} \\
\\ 
\texttt{Function signature:} \\
\texttt{(text: str) {-}> str} \\
\\ 
\texttt{You have 1 turns total.} \\
\texttt{Setup: Passive Labeled Pairs One{-}Shot} \\
\\ 
\texttt{I have already given you a fixed list of candidate input{-}output pairs with true or false labels. Solve using those statements only.} \\
\\ 
\texttt{You must make your best SOLVE attempt no later than the second{-}to{-}last turn.} \\
\texttt{Do not wait until the very end.} \\
\\ 
\texttt{You can use these actions:} \\
\texttt{{-} SOLVE: ```python ... ```} \\
\\ 
\texttt{Final answer format:} \\
\texttt{SOLVE: ```python} \\
\texttt{def solution(<value1>):} \\
\texttt{    ...} \\
\texttt{```} \\
\\ 
\texttt{Do not add import statements in your guess.} \\
\texttt{Anything you might need is already available, so just write the function body and return the answer.} \\
\\ 
\texttt{No further example requests are allowed in this one{-}shot baseline.} \\
\\ 
\texttt{I will reply with OUTPUT: lines only.} \\
\texttt{You must never write OUTPUT: yourself.} \\
\\ 
\texttt{Preloaded examples:} \\
\texttt{1. (('hello', 'hello!'), True)} \\
\texttt{2. (('aba', 'aba!\_alt'), False)} \\
\texttt{3. (('a{-}b', 'a{-}b!'), True)} \\
\texttt{4. (('hello world', 'hello world!\_alt'), False)} \\
\texttt{5. (('A', 'A!'), True)} \\
\texttt{6. (('a', 'a!\_alt'), False)} \\
\texttt{7. (('aA', 'aA!'), True)} \\
\texttt{8. (('', '!\_alt'), False)} \\
\texttt{9. (('   ', '   !'), True)} \\
\texttt{10. (('123', '123!\_alt'), False)} \\
\\ 
\texttt{Scoring:} \\
\texttt{{-} correctness is required for any non{-}zero score} \\
\texttt{{-} in interactive modes, solving with fewer information requests is better} \\
\\ 
\texttt{Input formatting reminders:} \\
\texttt{{-} strings must be quoted} \\
\texttt{{-} booleans must be True/False} \\
\texttt{{-} use commas between multiple values} \\
\texttt{{-} do not write any extra action tags in the same message} \\
\\ 
\texttt{Format example:} \\
\texttt{You write:} \\
\texttt{SOLVE: ```python} \\
\texttt{I answer:} \\
\texttt{def solution(...):} \\
\texttt{    ...} \\
\texttt{```} \\
\\ 
\texttt{Output only one valid action per turn.} \\
            }
        }
    }
     \\ \\

    \theutterance \stepcounter{utterance}  
    & \multicolumn{4}{p{0.6\linewidth}}{
        \cellcolor[rgb]{0.561,0.737,0.561}{
            \makecell[{{p{\linewidth}}}]{
                \texttt{\tiny{[P1$\rangle$GM]}}
                \texttt{def solution(text):} \\
\texttt{    \# Based on the provided examples, the pattern appears to be:} \\
\texttt{    \# {-} If the input string is non{-}empty, the output is the input string with an exclamation mark appended.} \\
\texttt{    \# {-} If the input string is empty, the output is not the input string with an exclamation mark (which would be just '!\_alt' in the example).} \\
\texttt{    \# {-} The label True corresponds to the case where the input string is non{-}empty.} \\
\texttt{    \# {-} The label False corresponds to the empty string case.} \\
\texttt{    \#} \\
\texttt{    \# Therefore, the function likely returns the input string with '!' appended if it's non{-}empty,} \\
\texttt{    \# otherwise returns a specific string (possibly '!\_alt' or similar).} \\
\texttt{    \#} \\
\texttt{    \# From the examples:} \\
\texttt{    \# ('hello', 'hello!') {-}> True} \\
\texttt{    \# ('', '!\_alt') {-}> False} \\
\texttt{    \# ('a', 'a!\_alt') {-}> False} \\
\texttt{    \# ('aA', 'aA!') {-}> True} \\
\texttt{    \# ('   ', '   !') {-}> True} \\
\texttt{    \#} \\
\texttt{    \# Wait, the last example ('   ', '   !') {-}> True suggests that any non{-}empty string (including spaces) results in appending '!'} \\
\texttt{    \#} \\
\texttt{    \# But the example ('hello', 'hello!') {-}> True, matches that.} \\
\texttt{    \# The example ('a', 'a!\_alt') {-}> False, indicates that if the string is 'a', the output is '} \\
            }
        }
    }
    & & \\ \\

    \theutterance \stepcounter{utterance}  
    & & & \multicolumn{2}{p{0.3\linewidth}}{
        \cellcolor[rgb]{0.9,0.9,0.9}{
            \makecell[{{p{\linewidth}}}]{
                \texttt{\tiny{[GM$|$GM]}}
                \texttt{Game Verdict: LOSS (turn limit reached)} \\
            }
        }
    }
    & & \\ \\

\end{supertabular}
}

\end{document}
